% Options for packages loaded elsewhere
\PassOptionsToPackage{unicode}{hyperref}
\PassOptionsToPackage{hyphens}{url}
\documentclass[
  ignorenonframetext,
]{beamer}
\newif\ifbibliography
\usepackage{pgfpages}
\setbeamertemplate{caption}[numbered]
\setbeamertemplate{caption label separator}{: }
\setbeamercolor{caption name}{fg=normal text.fg}
\beamertemplatenavigationsymbolsempty
% remove section numbering
\setbeamertemplate{part page}{
  \centering
  \begin{beamercolorbox}[sep=16pt,center]{part title}
    \usebeamerfont{part title}\insertpart\par
  \end{beamercolorbox}
}
\setbeamertemplate{section page}{
  \centering
  \begin{beamercolorbox}[sep=12pt,center]{section title}
    \usebeamerfont{section title}\insertsection\par
  \end{beamercolorbox}
}
\setbeamertemplate{subsection page}{
  \centering
  \begin{beamercolorbox}[sep=8pt,center]{subsection title}
    \usebeamerfont{subsection title}\insertsubsection\par
  \end{beamercolorbox}
}
% Prevent slide breaks in the middle of a paragraph
\widowpenalties 1 10000
\raggedbottom
\AtBeginPart{
  \frame{\partpage}
}
\AtBeginSection{
  \ifbibliography
  \else
    \frame{\sectionpage}
  \fi
}
\AtBeginSubsection{
  \frame{\subsectionpage}
}
\usepackage{iftex}
\ifPDFTeX
  \usepackage[T1]{fontenc}
  \usepackage[utf8]{inputenc}
  \usepackage{textcomp} % provide euro and other symbols
\else % if luatex or xetex
  \usepackage{unicode-math} % this also loads fontspec
  \defaultfontfeatures{Scale=MatchLowercase}
  \defaultfontfeatures[\rmfamily]{Ligatures=TeX,Scale=1}
\fi
\usepackage{lmodern}
\ifPDFTeX\else
  % xetex/luatex font selection
\fi
% Use upquote if available, for straight quotes in verbatim environments
\IfFileExists{upquote.sty}{\usepackage{upquote}}{}
\IfFileExists{microtype.sty}{% use microtype if available
  \usepackage[]{microtype}
  \UseMicrotypeSet[protrusion]{basicmath} % disable protrusion for tt fonts
}{}
\makeatletter
\@ifundefined{KOMAClassName}{% if non-KOMA class
  \IfFileExists{parskip.sty}{%
    \usepackage{parskip}
  }{% else
    \setlength{\parindent}{0pt}
    \setlength{\parskip}{6pt plus 2pt minus 1pt}}
}{% if KOMA class
  \KOMAoptions{parskip=half}}
\makeatother
\setlength{\emergencystretch}{3em} % prevent overfull lines
\providecommand{\tightlist}{%
  \setlength{\itemsep}{0pt}\setlength{\parskip}{0pt}}
\usepackage{bookmark}
\IfFileExists{xurl.sty}{\usepackage{xurl}}{} % add URL line breaks if available
\urlstyle{same}
\hypersetup{
  hidelinks,
  pdfcreator={LaTeX via pandoc}}

\author{\texorpdfstring{}{}}
\date{}

\begin{document}

\begin{frame}
\vspace*{2cm}

\begin{center}
\begin{minipage}{0.7\textwidth}
\centering
\Large\itshape
``The difference between theory and practice\\[0.3em]
is larger in practice than in theory.''
\vspace{1em}

\normalsize\upshape
--- Jan van de Snepscheut, Computer Scientist
\end{minipage}
\end{center}

\vspace{2cm}
\end{frame}

\begin{frame}{Abstract}
\protect\phantomsection\label{abstract}
This paper provides computational validation of the \textbf{Cognitive
Integrity Framework (CIF)} introduced in Part 1. We implement the
complete defense suite---cognitive firewalls, belief sandboxes, trust
calculus, and Byzantine-tolerant consensus---and evaluate performance
across production multiagent architectures.

\begin{block}{Contributions}
\protect\phantomsection\label{contributions}
\begin{itemize}
\tightlist
\item
  \textbf{Attack Corpus}: 950 cognitive attacks across four categories
  (prompt injection, trust exploitation, belief manipulation,
  coordination attacks)
\item
  \textbf{Cross-Architecture Validation}: Evaluation across Claude Code,
  AutoGPT, CrewAI, LangGraph, MetaGPT, and Camel
\item
  \textbf{Detection Performance}: 94\% detection rate with layered
  defenses; 20-25\% latency overhead
\item
  \textbf{Statistical Analysis}: Significance testing with large effect
  sizes (Cohen's d \textgreater{} 0.8), ablation studies, scalability
  benchmarks
\end{itemize}
\end{block}

\begin{block}{Key Findings}
\protect\phantomsection\label{key-findings}
\begin{enumerate}
\tightlist
\item
  \textbf{Composition Matters}: No individual defense achieves
  acceptable protection; layered composition yields multiplicative
  detection improvement
\item
  \textbf{Trust Decay Works}: Bounded delegation (δ\^{}d) prevents trust
  amplification across all architectures
\item
  \textbf{Architecture Vulnerability}: Peer-to-peer systems show largest
  relative improvement, confirming lateral movement analysis
\end{enumerate}

All notation follows definitions from Part 1 (Supplementary Section
S03).
\end{block}

\begin{block}{Paper Series}
\protect\phantomsection\label{paper-series}
\textbf{DOI}: 10.5281/zenodo.18364128

This is Part 2 of the \emph{Cognitive Security for Multiagent Operators}
series:

\begin{itemize}
\tightlist
\item
  \textbf{Part 1} (DOI: 10.5281/zenodo.18364119): Formal foundations and
  theoretical analysis
\item
  \textbf{Part 2} (this paper): Computational validation and
  implementation
\item
  \textbf{Part 3} (DOI: 10.5281/zenodo.18364130): Practical deployment
  guidance
\end{itemize}
\end{block}
\end{frame}

\end{document}
